\section{Introduction}
\label{sec:intro}

Language models increasingly translate questions into programs over structured data, making their
outputs inspectable and executable.  Parsers, schema linkers, and executors can expose malformed
syntax, unknown references, empty results, and runtime failures, but none establishes that a
program expresses the request.  This distinction matters during inference, when diagnostics can
guide revision, and during training, when false acceptance can turn semantic errors into targets.
We therefore ask how partial executable verification can provide feedback without obscuring its
limited authority.

Public procurement makes this boundary concrete.  Its records describe buyers, suppliers, awards,
values, dates, categories, and source notices.  An auditor may ask which buyers share suppliers
with a named authority, whether all awards in a period satisfy a condition, or which category has
the largest complete total.  These questions require joins across buyer and supplier roles, set operations,
aggregation, and access to the complete matching population in a declared data snapshot.
Retrieval alone cannot establish this population.  Let $R(q)$ contain the records that satisfy a
question $q$, and let $\widehat{R}_k(q)$ contain the top $k$ retrieved records.  Even when every
retrieved row is relevant, $\widehat{R}_k(q)\subset R(q)$ does not imply
$|\widehat{R}_k(q)|=|R(q)|$, nor does it preserve a sum, ranking, or set difference.  Executing a
generated program over the full snapshot avoids this finite retrieval limit, but it introduces the
separate problem of whether the program preserves the roles, constraints, and operation expressed
by the question.

We study this problem over a frozen graph of UK procurement records.  Open Contracting
Data Standard releases are reduced to a versioned snapshot, organisations are resolved
conservatively, and contract, organisation, category, and evidence objects are stored in a property
graph backed by Parquet.  Learned components propose a typed computation.  Deterministic
code grounds its fields and values, executes it over the matching population in the snapshot, and
applies evidence and release checks.  Deterministic execution fixes the result of a proposed
computation, but the model still chooses its roles, dependencies, and operation.  The design exposes
those choices to a defined set of checks.

Online checks have no access to the answer oracle.  A repairable failure may be returned to the
model for one revised proposal, while unsupported, ambiguous, and clean empty states terminate.
Offline curation adds a hidden typed oracle after execution.  Positive targets must pass the runtime
contract, satisfy the expected answer shape, and agree with that oracle, which contributes no
content to the student target.  The online revision changes no parameter or state across questions.

The full runtime contains two learned proposal stages.  A structured understander first predicts an
answer signature and a typed intent program.  Simple valid intents are compiled directly.  More
complex or vetoed intents are passed to a graph planner that exposes intermediate sets and their
dependencies.  Both routes enter the same grounding, execution, evidence, and abstention path.
The trace records grounded constraints, intermediate population sizes, and a bounded sample of
contract level provenance.  Figure \ref{fig:overview} summarises this control path.

\begin{figure}[t]
  \centering
  \IfFileExists{figs/fig01_system_overview.pdf}{%
    \includegraphics[width=\linewidth]{fig01_system_overview}%
  }{%
    \fbox{\parbox[c][1.45in][c]{0.94\linewidth}{\centering\small
      \textbf{Figure placeholder.} A procurement question runs from structured understanding to
      the selected compiler or planner branch, a typed variable graph, grounded intermediate sets,
      deterministic execution, release checks, and either an answer with evidence or abstention.
      The displayed values reconstruct a saved Step 1 artifact through deterministic replay on the
      current frozen backend.}}
  }
  \caption{Control path of the full procurement runtime.  The understander is always called.  A
  valid compilable intent takes the deterministic route, while other intents are passed to the
  typed graph planner.  Both branches are grounded and executed on the frozen graph before release
  checks.  Only a structured, repairable failure may trigger one revised proposal.  The displayed
  counts reconstruct a saved Step 1 artifact through deterministic replay rather than a stored full
  final test trace.}
  \label{fig:overview}
\end{figure}

The evaluation is organised around three questions.

\begin{itemize}
  \item \textbf{RQ1. Feedback authority.} Which properties can deterministic checks establish
  without access to the answer, and where do runtime acceptance, oracle agreement, and constraint
  retention diverge?
  \item \textbf{RQ2. Learning from checked executions.} How do stored Qwen and Llama checkpoints
  change across supervised fine tuning, rejection sampling fine tuning (RSFT), and preference
  optimisation,
  and do later stages improve monotonically?
  \item \textbf{RQ3. System effectiveness and portability.} How does the complete procurement
  runtime compare with direct generation, finite retrieval, cloud, and deterministic controls, and
  what planning competence transfers under a separate typed algebra protocol?
\end{itemize}

The main comparison uses the 2,025 questions outside the development subset containing 260
questions, which supports checkpoint analysis.  Targeted audits examine grounding and constraint
retention.  The Procurement Analytics Challenge Set (PACS) and WikiTableQuestions evaluate a
separate algebra planner using one model call without the full runtime.

The paper contributes a layered contract that separates learned proposals, deterministic online
checks, and offline oracle authority.  It also develops a training path that routes checked
candidates into direct targets, repairs, preferences, hard negatives, and abstention examples.  The
resulting systems are evaluated across two model families, targeted audits, and separate
composition and transfer protocols.  These protocols support different conclusions and are not
treated as points on one learning curve.
